\begin{center}
\textit{Opposite forms.}
\end{center}

\Article{159}
\textit{If the form $F$ contains the form $F'$, and this in turn contains the form $F''$, then the form $F$ also contains the form $F''$.}

Let the indeterminates of the forms $F,F',F''$ be $x,y$, $x',y'$, $x'',y''$ respectively, and suppose $F$ changes into $F'$ on putting
\[
x=\alpha x'+\beta y'.\quad y=\gamma x'+\delta y'
\]
and $F'$ into $F''$ on putting
\[
x'=\alpha'x''+\beta'y'',\quad y'=\gamma'x''+\delta'y''
\]
Clearly, $F$ will change into $F''$ on putting
\[
x=\alpha(\alpha'x''+\beta'y'')+\beta(\gamma'x''+\delta'y''),\quad
y=\gamma(\alpha'x''+\beta'y'')+\delta(\gamma'x''+\delta'y'')
\]
or
\[
x=(\alpha\alpha'+\beta\gamma')x''+(\alpha\beta'+\beta\delta')y'',\quad
y=(\gamma\alpha'+\delta\gamma')x''+(\gamma\beta'+\delta\delta')y''
\]
So $F$ contains $F''$.

\SourceLeaf{132}
Since
\[
(\alpha\alpha'+\beta\gamma')(\gamma\beta'+\delta\delta')
 -(\alpha\beta'+\beta\delta')(\gamma\alpha'+\delta\gamma')
 =(\alpha\delta-\beta\gamma)(\alpha'\delta'-\beta'\gamma')
\]
the result is positive if $\alpha\delta-\beta\gamma$ and $\alpha'\delta'-\beta'\gamma'$ are both positive or both negative, and negative if one is positive and the other negative. So the form $F$ contains the form $F''$ \textit{properly} if $F$ contains $F'$ in the same way as $F'$ contains $F''$, and \textit{improperly} if in different ways.

It follows that if we have any number of forms $F,F',F'',F'''$, etc., each containing the next, the first contains the last, \textit{properly} if the number of forms that contain their successor improperly is even, and \textit{improperly} if this number is odd.

\textit{If the form $F$ is equivalent to the form $F'$, and the form $F'$ to the form $F''$, then the form $F$ is equivalent to the form $F''$: properly if the form $F$ is equivalent to the form $F'$ in the same way as the form $F'$ is to the form $F''$, and improperly if in different ways.}

For since the forms $F,F'$ are respectively equivalent to $F',F''$, not only does each of the former contain the corresponding latter form, so that $F$ contains $F''$, but the latter also contain the former. So $F,F''$ are equivalent. The preceding result shows that $F$ contains $F''$ properly or improperly according as $F$ is equivalent to $F'$ in the same way as $F'$ is to $F''$, or in different ways; the same is true of $F''$ containing $F$. Thus, in the first case $F,F''$ are properly equivalent, and in the second, improperly equivalent.

\textit{The forms $(a,-b,c)$, $(c,b,a)$, $(c,-b,a)$ are equivalent to the form $(a,b,c)$: the first two improperly, and the last properly.}

For $axx+2bxy+cyy$ changes into $ax'x'-2bx'y'+cy'y'$ on putting $x=x'+0.y'$, $y=0.x'-y'$, and this transformation is improper because $1\times-1-0.0=-1$; into the form $cx'x'+2bx'y'+ay'y'$ by the improper transformation $x=0.x'+y'$, $y=x'+0.y'$; and into the form $cx'x'-2bx'y'+ay'y'$ by the proper transformation $x=0.x'-y'$, $y=x'+0.y'$.

It is therefore clear that any form equivalent to the form $(a,b,c)$ is \textit{properly} equivalent either to that form itself or to the form $(a,-b,c)$. Similarly, if any form contains the form $(a,b,c)$ or is contained in it, it \textit{properly} contains either the form $(a,b,c)$ or the form $(a,-b,c)$, or is \textit{properly} contained in one of them. We will call the forms $(a,b,c)$, $(a,-b,c)$ \textit{opposite}.

\SourceLeaf{133}
\begin{center}
\textit{Neighboring forms.}
\end{center}
\Article{160}
If the forms $(a,b,c)$, $(a',b',c')$ have the same determinant, and in addition $c=a'$ and $b\equiv-b'\; (\mathrm{mod.}\,c)$, or $b+b'\equiv0\; (\mathrm{mod.}\,c)$, we will call them \textit{neighboring} forms. When a more precise description is needed, we will call the first a neighbor of the second \textit{on the left}, and the second a neighbor of the first \textit{on the right}.

For \textit{example}, the form $(7,3,2)$ is a neighbor of the form $(3,4,7)$ on the right, and the form $(3.1.3)$ is a neighbor of its opposite $(3,-1,3)$ on both sides.

\textit{Neighboring forms are always properly equivalent.} For the form $axx+2bxy+cyy$ changes into the neighboring form $cx'x'+2b'x'y'+c'y'y'$ under the substitution $x=-y'$, $y=x'+\frac{b+b'}{c}y'$ (which is proper because $0\times\frac{b+b'}{c}-1\times-1=1$), as is easily proved by expansion using the equation $bb-ac=b'b'-cc'$; and $\frac{b+b'}{c}$ is an integer by hypothesis. --- These definitions and conclusions do not apply if $c=a'=0$. This case can occur only for forms whose determinant is a square.

The forms $(a,b,c)$, $(a',b',c')$ are properly equivalent if $a=a'.\ b\equiv b'\; (\mathrm{mod.}\,a)$. For the form $(a,b,c)$ is properly equivalent to the form $(c,-b,a)$ (the preceding article), and the latter is a neighbor of the form $(a',b',c')$ on the left.

\begin{center}
\textit{Common divisors of the coefficients of forms.}
\end{center}
\Article{161}
\textit{If the form $(a,b,c)$ contains the form $(a',b',c')$, every common divisor of the numbers $a,b,c$ also divides the numbers $a',b',c'$, and every common divisor of the numbers $a,2b,c$ divides $a',2b',c'$.}

For if the form $axx+2bxy+cyy$ changes under the substitutions $x=\alpha x'+\beta y'$, $y=\gamma x'+\delta y'$ into the form $a'x'x'+2b'x'y'+c'y'y'$, we have the equations
\[
\begin{aligned}
a\alpha\alpha+2b\alpha\gamma+c\gamma\gamma&=a'\\
a\alpha\beta+b(\alpha\delta+\beta\gamma)+c\gamma\delta&=b'\\
a\beta\beta+2b\beta\delta+c\delta\delta&=c'
\end{aligned}
\]
from which the proposition follows immediately (for its second part, use $2a\alpha\beta+2b(\alpha\delta+\beta\gamma)+2c\gamma\delta=2b'$ in place of the second equation).

\SourceLeaf{134}
It follows that the greatest common divisor of the numbers $a,b(2b),c$ also divides the greatest common divisor of the numbers $a',b'(2b'),c'$. If, in addition, the form $(a',b',c')$ contains the form $(a,b,c)$, that is, if the forms are equivalent, the greatest common divisor of the numbers $a,b(2b),c$ equals the greatest common divisor of the numbers $a',b'(2b'),c'$, since each must then divide the other. Thus, if in this case $a,b(2b),c$ have no common divisor, that is, if their greatest common divisor $=1$, then $a',b'(2b'),c'$ also have no common divisor.

\Needspace{18\baselineskip}
\begin{center}
\textit{Connection between all similar transformations of a given form into a given form.}
\end{center}
\Article{162}
\textsc{Problem.} \textit{If the form}
\[
AXX+2BXY+CYY\quad\mathinner{.\,.\,.}\quad F
\]
\textit{contains the form}
\[
axx+2bxy+cyy\quad\mathinner{.\,.\,.\,.\,.}\quad f
\]
\textit{and one transformation of the former into the latter is given, derive from it all the other transformations similar to it.}

\textit{Solution.} Let the given transformation be $X=\alpha x+\beta y$, $Y=\gamma x+\delta y$. First suppose that another similar transformation $X=\alpha'x+\beta'y$, $Y=\gamma'x+\delta'y$ is given, and investigate what follows from this. For the forms and their determinants, write $F,f,\ =D,d$, and put $\alpha\delta-\beta\gamma=e$, $\alpha'\delta'-\beta'\gamma'=e'$. We have (Article 157) $d=Dee=De'e'$. Since by hypothesis $e,e'$ have the same sign, $e=e'$. We have the following six equations:
\[
\begin{aligned}
A\alpha\alpha+2B\alpha\gamma+C\gamma\gamma&=a&\quad\mathinner{.\,.\,.\,.\,.\,.\,.}\quad&[1]\\
A\alpha'\alpha'+2B\alpha'\gamma'+C\gamma'\gamma'&=a&\quad\mathinner{.\,.\,.\,.\,.\,.\,.}\quad&[2]\\
A\alpha\beta+B(\alpha\delta+\beta\gamma)+C\gamma\delta&=b&\quad\mathinner{.\,.\,.\,.\,.\,.}\quad&[3]\\
A\alpha'\beta'+B(\alpha'\delta'+\beta'\gamma')+C\gamma'\delta'&=b&\quad\mathinner{.\,.\,.\,.\,.\,.}\quad&[4]\\
A\beta\beta+2B\beta\delta+C\delta\delta&=c&\quad\mathinner{.\,.\,.\,.\,.\,.\,.}\quad&[5]\\
A\beta'\beta'+2B\beta'\delta'+C\delta'\delta'&=c&\quad\mathinner{.\,.\,.\,.\,.\,.\,.}\quad&[6]
\end{aligned}
\]
If, for brevity, we denote the numbers
\[
\begin{gathered}
A\alpha\alpha'+B(\alpha\gamma'+\gamma\alpha')+C\gamma\gamma'\\
A(\alpha\beta'+\beta\alpha')+B(\alpha\delta'+\beta\gamma'+\gamma\beta'+\delta\alpha')+C(\gamma\delta'+\delta\gamma')\\
A\beta\beta'+B(\beta\delta'+\delta\beta')+C\delta\delta'
\end{gathered}
\]
\SourceLeaf{135}
by $a',2b',c'$, we derive the following new equations from the preceding ones\HistoricalNote{*)}{These equations arise as follows: $7$ comes from $1.2$ (that is, by multiplying equation (1) by equation (2), or rather by multiplying the left side of the former by the left side of the latter, and its right side by the right side of the latter, and equating the products); $8$ from $1.4+2.3$; the next, unnumbered equation from $1.6+2.5+3.4+3.4$; the next unnumbered one from $3.4$; $11$ from $3.6+4.5$; $12$ from $5.6$. We will always use similar notation in what follows. We must leave the expansion to the reader.}:
\[
\begin{aligned}
a'a'-D(\alpha\gamma'-\gamma\alpha')^2&=aa&\quad\mathinner{.\,.\,.\,.\,.\,.\,.\,.}\quad&[7]\\
2a'b'-D(\alpha\gamma'-\gamma\alpha')(\alpha\delta'+\beta\gamma'-\gamma\beta'-\delta\alpha')&=2ab&\quad\mathinner{.\,.}\quad&[8]\\
4b'b'-D\bigl((\alpha\delta'+\beta\gamma'-\gamma\beta'-\delta\alpha')^2+2ee'\bigr)&=2bb+2ac
\end{aligned}
\]
Adding $2Dee'=2d=2bb-2ac$ gives
\[
\begin{aligned}
4b'b'-D(\alpha\delta'+\beta\gamma'-\gamma\beta'-\delta\alpha')^2&=4bb&\quad\mathinner{.\,.\,.\,.\,.}\quad&[9]\\
a'c'-D(\alpha\delta'-\gamma\beta')(\beta\gamma'-\delta\alpha')&=bb
\end{aligned}
\]
and subtracting $D(\alpha\delta-\beta\gamma)(\alpha'\delta'-\beta'\gamma')=bb-ac$ gives
\[
\begin{aligned}
a'c'-D(\alpha\gamma'-\gamma\alpha')(\beta\delta'-\delta\beta')&=ac&\quad\mathinner{.\,.\,.\,.\,.\,.}\quad&[10]\\
2b'c'-D(\alpha\delta'+\beta\gamma'-\gamma\beta'-\delta\alpha')(\beta\delta'-\delta\beta')&=2bc&\quad\mathinner{.\,.}\quad&[11]\\
c'c'-D(\beta\delta'-\delta\beta')^2&=cc&\quad\mathinner{.\,.\,.\,.\,.\,.\,.\,.}\quad&[12]
\end{aligned}
\]
Now let the greatest common divisor of the numbers $a,2b,c$ be $m$, and choose the numbers $\mathfrak{A},\mathfrak{B},\mathfrak{C}$ so that
\[
\mathfrak{A}a+2\mathfrak{B}b+\mathfrak{C}c=m
\]
(Article 40). Multiply the equations $7,8,9,10,11,12$ by $\mathfrak{A}\mathfrak{A},\allowbreak 2\mathfrak{A}\mathfrak{B},\allowbreak \mathfrak{B}\mathfrak{B},\allowbreak 2\mathfrak{A}\mathfrak{C},\allowbreak 2\mathfrak{B}\mathfrak{C},\allowbreak \mathfrak{C}\mathfrak{C}$ respectively and add the products. If, for brevity, we put
\[
\begin{aligned}
\mathfrak{A}a'+2\mathfrak{B}b'+\mathfrak{C}c'&=T&\quad\mathinner{.\,.\,.\,.\,.\,.\,.\,.}\quad&[13]\\
\mathfrak{A}(\alpha\gamma'-\gamma\alpha')&+\mathfrak{B}(\alpha\delta'+\beta\gamma'-\gamma\beta'-\delta\alpha')&&\\
{}+\mathfrak{C}(\beta\delta'-\delta\beta')&=U&\quad\mathinner{.\,.}\quad&[14]
\end{aligned}
\]
where $T,U$ are clearly integers, we obtain
\[
TT-DUU=mm
\]
So we have reached this elegant conclusion: \textit{any two similar transformations of the form $F$ into $f$ give an integer solution of the indeterminate equation $tt-Duu=mm$}, namely $t=T,\ u=U$. Since in
\SourceLeaf{136}
our reasoning we did not assume the transformations to be \textit{different}, even \textit{one} transformation considered twice must give a solution. In that case, however, because $\alpha'=\alpha,\ \beta'=\beta$, etc., we have $a'=a,\ b'=b,\ c'=c$, and so $T=m,\ U=0$, a solution that is obvious in itself.

Now consider the first transformation and the solution of the indeterminate equation as known, and investigate how the second transformation can be derived from them, or how $\alpha',\beta',\gamma',\delta'$ depend on $\alpha,\beta,\gamma,\delta,\ T,U$. For this purpose, first multiply equation $[1]$ by $\delta\alpha'-\beta\gamma'$, $[2]$ by $\alpha\delta'-\gamma\beta'$, $[3]$ by $\alpha\gamma'-\gamma\alpha'$, $[4]$ by $\gamma\alpha'-\alpha\gamma'$, and add the products. This gives
\[
(e+e')a'=(\alpha\delta'-\beta\gamma'-\gamma\beta'+\delta\alpha')a\quad\mathinner{.\,.\,.\,.\,.\,.}\quad[15]
\]
Similarly, from
\[
(\delta\beta'-\beta\delta')([1]-[2])+(\alpha\delta'-\beta\gamma'-\gamma\beta'+\delta\alpha')([3]+[4])+(\alpha\gamma'-\gamma\alpha')([5]-[6])
\]
we get
\[
2(e+e')b'=2(\alpha\delta'-\beta\gamma'-\gamma\beta'+\delta\alpha')b\quad\mathinner{.\,.\,.\,.\,.}\quad[16]
\]
Finally, from $(\delta\beta'-\beta\delta')([3]-[4])+(\alpha\delta'-\gamma\beta')[5]+(\delta\alpha'-\beta\gamma')[6]$ we get
\[
(e+e')c'=(\alpha\delta'-\beta\gamma'-\gamma\beta'+\delta\alpha')c\quad\mathinner{.\,.\,.\,.\,.\,.}\quad[17]
\]
Substituting these values $(15,16,17)$ into $13$ gives
\[
(e+e')T=(\alpha\delta'-\beta\gamma'-\gamma\beta'+\delta\alpha')(\mathfrak{A}a+2\mathfrak{B}b+\mathfrak{C}c)
\]
or
\[
2eT=(\alpha\delta'-\beta\gamma'-\gamma\beta'+\delta\alpha')m\quad\mathinner{.\,.\,.\,.\,.\,.\,.}\quad[18]
\]
from which $T$ can be derived much more easily than from $[13]$. --- Combining this equation with $15,16,17$ gives $ma'=Ta,\ 2mb'=2Tb,\ mc'=Tc$. Substituting these values of $a',2b',c'$ into equations $7\text{---}12$ and writing $TT$ as $mm+DUU$, we obtain, after the required changes,
\[
\begin{aligned}
(\alpha\gamma'-\gamma\alpha')^2mm&=aaUU\\
(\alpha\gamma'-\gamma\alpha')(\alpha\delta'+\beta\gamma'-\gamma\beta'-\delta\alpha')mm&=2abUU
\end{aligned}
\]
\SourceLeaf{137}
\[
\begin{aligned}
(\alpha\delta'+\beta\gamma'-\gamma\beta'-\delta\alpha')^2mm&=4bbUU\\
(\alpha\gamma'-\gamma\alpha')(\beta\delta'-\delta\beta')mm&=acUU\\
(\alpha\delta'+\beta\gamma'-\gamma\beta'-\delta\alpha')(\beta\delta'-\delta\beta')mm&=2bcUU\\
(\beta\delta'-\delta\beta')^2mm&=ccUU
\end{aligned}
\]
Using equation $[14]$ and the equation $\mathfrak{A}a+2\mathfrak{B}b+\mathfrak{C}c=m$, we easily derive the following (multiply the first, second and fourth; the second, third and fifth; the fourth, fifth and sixth equations by $\mathfrak{A},\mathfrak{B},\mathfrak{C}$ respectively, and add the products):
\[
\begin{aligned}
(\alpha\gamma'-\gamma\alpha')Umm&=maUU\\
(\alpha\delta'+\beta\gamma'-\gamma\beta'-\delta\alpha')Umm&=2mbUU\\
(\beta\delta'-\delta\beta')Umm&=mcUU
\end{aligned}
\]
and then, dividing by $mU$\HistoricalNote{*)}{This would not be allowed if $U=0$. In that case, however, equations $19,20,21$ would follow immediately from the first, third and sixth of the preceding equations.},
\[
\begin{aligned}
aU&=(\alpha\gamma'-\gamma\alpha')m&\quad\mathinner{.\,.\,.\,.\,.\,.\,.\,.\,.}\quad&[19]\\
2bU&=(\alpha\delta'+\beta\gamma'-\gamma\beta'-\delta\alpha')m&\quad\mathinner{.\,.\,.\,.\,.}\quad&[20]\\
cU&=(\beta\delta'-\delta\beta')m&\quad\mathinner{.\,.\,.\,.\,.\,.\,.\,.\,.}\quad&[21]
\end{aligned}
\]
Any of these equations gives $U$ much more easily than $[14]$ does. --- It also follows that however $\mathfrak{A},\mathfrak{B},\mathfrak{C}$ are chosen (which can be done in infinitely many ways), both $T$ and $U$ retain their respective values.

Now multiplying equation $18$ by $\alpha$, $19$ by $2\beta$, and $20$ by $-\alpha$, and adding gives
\[
2\alpha eT+2(\beta a-\alpha b)U=2(\alpha\delta-\beta\gamma)\alpha'm=2e\alpha'm
\]
Similarly, from $\beta[18]+\beta[20]-2\alpha[21]$ we get
\[
2\beta eT+2(\beta b-\alpha c)U=2(\alpha\delta-\beta\gamma)\beta'm=2e\beta'm
\]
Further, from $\gamma[18]+2\delta[19]-\gamma[20]$ we get
\[
2\gamma eT+2(\delta a-\gamma b)U=2(\alpha\delta-\beta\gamma)\gamma'm=2e\gamma'm
\]
Finally, from $\delta[18]+\delta[20]-2\gamma[21]$ we get
\[
2\delta eT+2(\delta b-\gamma c)U=2(\alpha\delta-\beta\gamma)\delta'm=2e\delta'm
\]

\SourceLeaf{138}
Substituting in these formulas the values of $a,b,c$ from $1,3,5$ gives
\[
\begin{aligned}
\alpha'm&=\alpha T-(\alpha B+\gamma C)U\\
\beta'm&=\beta T-(\beta B+\delta C)U\\
\gamma'm&=\gamma T+(\alpha A+\gamma B)U\\
\delta'm&=\delta T+(\beta A+\delta B)U\text{\textsuperscript{*)}}
\end{aligned}
\]
\HistoricalDisplayNote{*)}{From this we easily derive
\[
\begin{aligned}
AeU&=(\delta\gamma'-\gamma\delta')m\\
2BeU&=(\alpha\delta'-\delta\alpha'+\gamma\beta'-\beta\gamma')m\\
CeU&=(\beta\alpha'-\alpha\beta')m
\end{aligned}
\]}

The preceding analysis shows that there is no transformation of the form $F$ into $f$ similar to the given one that is not included in the formula
\[
\begin{aligned}
X&=\tfrac{1}{m}(\alpha t-(\alpha B+\gamma C)u)x+\tfrac{1}{m}(\beta t-(\beta B+\delta C)u)y\\
Y&=\tfrac{1}{m}(\gamma t+(\alpha A+\gamma B)u)x+\tfrac{1}{m}(\delta t+(\beta A+\delta B)u)y\quad\mathinner{.\,.\,.}\quad(\mathrm{I})
\end{aligned}
\]
where $t,u$ stand for all integers satisfying the equation $tt-Duu=mm$. We cannot yet conclude that all values of $t,u$ satisfying that equation give suitable transformations when substituted into formula $(\mathrm{I})$. But:

1. The form $F$ always changes, under a substitution coming from any such values of $t,u$, into the form $f$. This is easily confirmed by expansion using equations $1,3,5$ and the equation $tt-Duu=mm$. We omit the calculation for brevity; it is longer than it is difficult.

2. Every transformation derived from the formula is similar to the given one. For
\[
\begin{aligned}
&\tfrac{1}{m}(\alpha t-(\alpha B+\gamma C)u)\times\tfrac{1}{m}(\delta t+(\beta A+\delta B)u)\\
&\quad-\tfrac{1}{m}(\beta t-(\beta B+\delta C)u)\times\tfrac{1}{m}(\gamma t+(\alpha A+\gamma B)u)\\
&\qquad=\tfrac{1}{mm}(\alpha\delta-\beta\gamma)(tt-Duu)=\alpha\delta-\beta\gamma
\end{aligned}
\]
3. If the forms $F,f$ have unequal determinants, formula $(\mathrm{I})$ may for some values of $t,u$ give substitutions involving \textit{fractions}, which must therefore be rejected. All the others will be suitable transformations, and there will be no others beyond them.

4. If, however, the forms $F,f$ have the same determinant and so are \textit{equivalent}, formula $(\mathrm{I})$ will give no transformations involving fractions
\SourceLeaf{139}
and will therefore give the complete solution of the problem in this case. We prove this as follows.

The theorem of the preceding article shows that in this case $m$ is also a common divisor of the numbers $A,2B,C$. Since $tt-Duu=mm$, we have $tt-BBuu=mm-ACuu$, so $tt-BBuu$ is divisible by $mm$. Then certainly $4tt-4BBuu$ is also divisible by it, and (because $2B$ is divisible by $m$) so is $4tt$ by $mm$, and therefore $2t$ by $m$. Thus $\frac{2}{m}(t+Bu),\ \frac{2}{m}(t-Bu)$ are integers, and, since their difference $\frac{4}{m}Bu$ is even, they are either both even or both odd. If both were odd, their product would also be odd. But as four times the number $\frac{1}{mm}(tt-BBuu)$, which we have just shown to be an integer, this product is necessarily even. So that case is impossible, and $\frac{2}{m}(t+Bu),\ \frac{2}{m}(t-Bu)$ are always even, giving integers $\frac{1}{m}(t+Bu),\ \frac{1}{m}(t-Bu)$. It follows without difficulty that all four coefficients in $(\mathrm{I})$ are always integers. \mbox{\textit{Q.E.D.}}

It follows that once all solutions of the equation $tt-Duu=mm$ are known, all transformations of the form $(A,B,C)$ into $(a,b,c)$ similar to the given transformation can be derived from them. We will show later how to find those solutions. Here we note only that their number is always finite when $D$ is negative, or positive and a square; it is infinite when $D$ is positive and not a square. In the latter case, if also the equality $D$$=d$ does not hold (point $3^0$ above), one would also have to investigate how to distinguish in advance the values of $t,u$ that give substitutions free of fractions from those that give fractional substitutions. But for this case we will give another method below, free of this drawback (Article 214).

\textit{Example.} The form $xx+2yy$ changes under the proper substitution $x=2x'+7y'$, $y=x'+5y'$ into the form $(6,24,99)$. We want \textit{all} proper transformations of the former form into the latter. Here $D=-2,\ m=3$, so the equation to solve is $tt+2uu=9$. It is satisfied in six different ways, namely by putting $t=3,-3,1,-1,1,-1;\ u=0,0,2,2,-2,-2$ respectively. The third and sixth solutions give fractional substitutions and must therefore be rejected. The others give four substitutions:
\[
x=\left\{\begin{aligned}
&2x'+7y'\\
&-2x'-7y'\\
&-2x'-9y'\\
&2x'+9y'
\end{aligned}\right.
\qquad
y=\left\{\begin{aligned}
&x'+5y'\\
&-x'-5y'\\
&x'+3y'\\
&-x'-3y'
\end{aligned}\right.
\]
(the first of which is the given one).
